\section{A second moment for the latent rows}
\label{sec:moments}

The independent-root comparison leaves the latent kernel
$\mathcal L$ of \eqref{eq:20}.  Its conditional first moments are bounded
uniformly over every site type by \eqref{eq:23}.  For the sampling argument
we also need an integrated second-moment estimate.  A small bound on each
individual representation does not suffice, because an edge may have
several representations. Expanding the square and weighting one
representation will identify the measure under which that multiplicity
must be controlled.

The grid size $L$ and tolerance $\tau$ will be chosen once below,
consistently with Section~\ref{sec:arithmetic}. In every limit
$B\to\infty$ in this section, those choices, $\eta>0$, the block constant
$C_6$, $C_*$, and the smooth weights are held fixed.
All estimates are uniform in the positions and in the smooth parameter
$s$ on the fixed compact range used in Section~\ref{sec:graph}.
Constants may depend on the fixed parameters but not on the particular
first representation.

\begin{proposition}[Integrated row moments]
\label{prop:moments:rows}
For a sufficiently fine fixed regularity grid and then a sufficiently
small fixed tolerance, the independent-site latent kernel satisfies
\begin{equation}
 \sum_{k\ne i}\E\,\mathcal L_{ik}^{2}\ll B^{-.21}
 \qquad (1\le i\le M).
 \label{eq:24}
\end{equation}
Moreover, its normalized total mass has bounded second moment:
\[
 \E\left(\frac1M\sum_{i,k}\mathcal L_{ik}\right)^2\ll1.
\]
For every set $I\subseteq\{1,\ldots,M\}$, the same estimates hold after
replacing $\mathcal L_{ik}$ by
$\ind_{\{i,k\in I\}}\mathcal L_{ik}$; the normalization remains $M$.
\end{proposition}

Write $h_\tau=h_0+\tau\log2$.  There are five truncated weights in
an individual summand of \eqref{eq:20}.  Since $k_B\ll1$ and the smooth
factors are bounded, \eqref{eq:8} gives the uniform bound
\begin{equation}
\begin{gathered}
 \text{one contribution to }\mathcal L_{ik}
 \ \le B^{-\kappa_1+5\tau\log2+o(1)},
 \\
 \kappa_1=1-5h_0=0.232867951\ldots> .2328.
 \label{eq:25}
\end{gathered}
\end{equation}
\subsection{The measure obtained from a first representation}

Fix an allowed pair $i,k=i+j$. Let
$\mathcal C_j$ be the finite set of positive triples $(a,b,c)$ satisfying
\[
 a-b=jc,\qquad
 B\le\log c\le2B,\qquad Tc\le a,b\le2Tc,
\]
whose entries are pairwise coprime, rough, squarefree, and regular.
These are precisely the coefficient conditions for a candidate; the
two remainder regularity conditions still depend on the site types.
We initially take $j>0$; the negative case follows by exchanging the
endpoints. Fix $(a,b,c)\in\mathcal C_j$. For an integer $v$ whose prime
factors belong to $\mathcal P$, write $\mathcal P(v)$ for its set of prime
factors.

In the independent-site model, the event
$\mathcal P(b)\subset S_i$, $\mathcal P(a)\subset S_k$ has probability
$1/(ab)$.  Conditional on that event, put
\[
 \mathcal T_b=S_i\setminus\mathcal P(b),
 \qquad
 \mathcal T_a=S_k\setminus\mathcal P(a).
\]
These sets are independent, and at each prime not excluded by the
corresponding coefficient the inclusion probability is $1/p$.
Weighting one such inclusion by its factor $1/2$ in $K_0$ changes this
probability to
\[
 \frac{(1/p)/2}{1-1/(2p)}=\frac1{2p-1}.
\]
Thus weighting by $K_0(\mathcal T_b)K_0(\mathcal T_a)$ gives a product
probability measure, denoted by $\Prob^{\mathrm{tilt}}_{a,b}$, under which
the two remainder sets are independent and their available prime
indicators have probabilities $1/(2p-1)$.

For a coefficient $v$ the normalizing factor is
\[
 \Xi(v)=R^{1/2}
       \prod_{p\in\mathcal P\setminus\mathcal P(v)}
                    \left(1-\frac1{2p}\right).
\]
The full product with no exclusions is bounded by Mertens' estimate.
Furthermore $\log v=O(B)$ and every prime of $v$ exceeds $P_0$, so
\[
 \sum_{p\mid v}\frac1p\ll\frac{B}{P_0}.
\]
Removing these factors changes the product by $1+o(1)$.  Consequently
$\Xi(v)\ll1$ uniformly for the coefficients under consideration.
For every nonnegative function $H$ of the two remainder sets we therefore
have the exact change-of-measure identity
\[
 \begin{split}
 &\E\big[
   \ind_{\mathcal P(b)\subset S_i,\,\mathcal P(a)\subset S_k}
   K_0(\mathcal T_b)K_0(\mathcal T_a)H(\mathcal T_b,\mathcal T_a)
       \big]\\
 &\hspace{25mm}
   =\frac{\Xi(b)\Xi(a)}{ab}
       \E^{\mathrm{tilt}}_{a,b}H(\mathcal T_b,\mathcal T_a).
 \end{split}
\]
The indicators that the two first remainders are regular will stay
inside $H$.  In particular, we never divide by their probability.

Given these remainder sets, reconstruct
$S_i=\mathcal P(b)\cup\mathcal T_b$ and
$S_k=\mathcal P(a)\cup\mathcal T_a$, and let
$N_{a,b,c}(\mathcal T_b,\mathcal T_a)$ be the number of valid candidates
at this pair for the reconstructed types, including the first triple
whenever it is valid.

We can now see exactly how this measure enters the square.
Expand the square of the sum of nonnegative candidate contributions,
designate one candidate as the first representation, and bound every
alternative contribution by \eqref{eq:25}. The first contribution retains
the two factors $K_0(\mathcal T_b)K_0(\mathcal T_a)$ and their regularity
indicators. The change-of-measure identity gives
\begin{equation}
\begin{split}
 \E\,\mathcal L_{ik}^{2}
 &\ll B^{-\kappa_1+5\tau\log2+o(1)}
 \frac1B\sum_{(a,b,c)\in\mathcal C_j}
       \frac{A(a)A(b)A(c)}{ab}\,\Xi(a)\Xi(b)\\
 &\qquad\cdot
 \E^{\mathrm{tilt}}_{a,b}
 \big[
   \ind_{\mathcal T_b\ \mathrm{regular}}
   \ind_{\mathcal T_a\ \mathrm{regular}}
   N_{a,b,c}(\mathcal T_b,\mathcal T_a)
 \big].
\end{split}
\label{eq:moments-first-representation}
\end{equation}
The implied constant absorbs $k_B$ and the bounded smooth factors.
Thus the remaining multiplicity question is the following uniform
estimate; the coefficient sum will then be bounded by \eqref{eq:7}.

\begin{lemma}[Weighted representation multiplicity]
\label{lem:moments:multiplicity}
The grid and tolerance can be chosen so that, uniformly for
$(a,b,c)\in\mathcal C_j$ and $\eta T\le |j|\le T$,
\[
 \E^{\mathrm{tilt}}_{a,b}
   \big[
    \ind_{\mathcal T_b\ \mathrm{regular}}
    \ind_{\mathcal T_a\ \mathrm{regular}}
    N_{a,b,c}(\mathcal T_b,\mathcal T_a)
   \big]
 \ll B^{.009}.
\]
\end{lemma}

\begin{proof}
Every alternative candidate has a unique decomposition
\[
 b'=d_1y,\quad d_1\mid b,\quad\mathcal P(y)\subset\mathcal T_b,
 \qquad
 a'=e_1z,\quad e_1\mid a,\quad\mathcal P(z)\subset\mathcal T_a,
 \qquad a'-b'=jc'.
\]
Here all the products are squarefree, and $y$ is coprime to $b$ while
$z$ is coprime to $a$.  The coefficient $c'$ is determined once the
four products are chosen.  We may discard any of the alternative
candidate conditions for upper bounds; however, we retain the
coefficient regularity conditions that are used below.  Since
$a',b'\le2T\exp(2B)$, we have $\log y,\log z\le3B$ for all sufficiently
large $B$.

Set $\Delta=1/L$, $g_0=3\Delta$, and
$U=\max(\log y,\log z)$. We count the alternatives with
$U\le B^{g_0}$ directly from prefix regularity. When $U>B^{g_0}$,
fixing the omissions and the new product at one endpoint places the
other new product in one residue class modulo $|j|$, through
$e_1z-d_1y=jc'$. The resulting progression density is the gain that
will pay for the remaining subset choices.

For a grid point $g<1$, define
\[
 \mathcal Q_g=\{p\in\mathcal P:\log p\le B^g\},
 \qquad \mathcal Q_1=\mathcal P,
\]
and write $\omega_g(v)=|\mathcal P(v)\cap\mathcal Q_g|$, with the
same notation for a prime set.  Regularity gives
\[
 (g/2-\tau)\ell\le\omega_g(v)\le(g/2+\tau)\ell
\]
for each of the first and alternative coefficients, and for each
regular first remainder.  Mertens' estimate, on this fixed grid, gives
\[
 \sum_{p\in\mathcal Q_g}\frac1p=g\ell+o(\ell).
\]
The $o(\ell)$ is uniform over the finitely many grid points.

\paragraph{Omissions above an addition prefix.}
Suppose that all primes of $z$ and $y$ lie in $\mathcal Q_g$.
Let $D_a=\mathcal P(a)\setminus\mathcal P(e_1)$ be the primes omitted
from $a$.  Regularity of the total counts of $a,a'$ gives
\[
 \bigl||D_a|-\omega(z)\bigr|\le2\tau\ell.
\]
Regularity at $g$ gives the analogous inequality
\[
 \bigl||D_a\cap\mathcal Q_g|-\omega(z)\bigr|\le2\tau\ell.
\]
Subtracting shows that at most $4\tau\ell$ primes above the prefix
can be omitted.  The same conclusion holds for the omissions from $b$.
When $g=1$ there are no primes above the prefix.

Write
\[
 H(u)=-u\log u-(1-u)\log(1-u),\qquad 0\le u\le1,
\]
with the endpoint values defined by continuity.
For completeness, if $m\le8\tau\ell$, the number of such subsets is
at most $2^m\le2^{8\tau\ell}$.  For
$8\tau\ell<m\le(1/2+\tau)\ell$, the binomial bound
\[
 \sum_{q\le4\tau\ell}\binom mq
 \le (m+1)\exp\bigl(mH(4\tau\ell/m)\bigr)
 \le (m+1)\exp\left((1/2+\tau)\ell
                 H\left(\frac{4\tau}{1/2+\tau}\right)\right)
\]
applies: $H$ is increasing up to $1/2$, and $mH(k/m)$ is increasing
in $m$ for fixed $0<k<m$.  Thus the choices at both coefficients
together cost $B^{\varepsilon_{\mathrm{high}}(\tau)+o(1)}$, where
$\varepsilon_{\mathrm{high}}(\tau)\to0$ as $\tau\to0$.
The estimate is uniform also when the relevant set is empty.

\paragraph{Small additions.}
If
\[
 U\le B^{g_0},
\]
all additions lie in $\mathcal Q_{g_0}$.  On the two first-remainder
regularity events, the four relevant prefix sets---the prime sets of
$a,b$ and the two first remainders---have at most
$(2g_0+4\tau)\ell$ primes altogether.  Allowing arbitrary subset choices
there, and then the higher omissions just bounded, gives the
pointwise estimate
\[
 N_{\mathrm{small}}
 \le 2^{(2g_0+4\tau)\ell}
       B^{\varepsilon_{\mathrm{high}}(\tau)+o(1)}
 =B^{(2g_0+4\tau)\log2+
                  \varepsilon_{\mathrm{high}}(\tau)+o(1)}.
\]
This also bounds its expectation with the first-remainder regularity
indicators.  By taking $L$ large and then $\tau$ small, its exponent
apart from $o(1)$ can be made strictly smaller than $.008$.

\paragraph{Large-addition classes and combinatorial choices.}
For $U>B^{g_0}$ use the classes
\[
 B^{g-\Delta}<U\le B^g
       \quad(g_0<g<1),
 \qquad
 B^{1-\Delta}<U\le3B\quad(g=1),
\]
where $g$ ranges over grid points.  All addition primes again belong to
$\mathcal Q_g$.  It suffices to treat the portion of a class in which
$\log z>B^{g-\Delta}$.  The portion with $\log y>B^{g-\Delta}$ is
estimated by the same argument with the endpoints exchanged; the
numeric sieve below permits either sign for its second slope.  If both
inequalities hold, counting the alternative twice only increases the
upper bound.

Define the exact normalized counts and their clipped values by
\[
 r_z^* =\frac{\omega(z)}{g\ell},\qquad
 r_y^* =\frac{\omega(y)}{g\ell},\qquad
 r_z=\min(r_z^*,1/2),\qquad r_y=\min(r_y^*,1/2).
\]
On the first-remainder regularity events,
$0\le r_z^*,r_y^*\le1/2+\tau/g$.  There are $O(\ell^2)$ possible pairs of
integer counts.  We estimate one such pair at a time.

The number of omitted primes in the $a$ prefix is
$r_z^*g\ell+O(\tau\ell)$, among
$g\ell/2+O(\tau\ell)$ available primes.  Applying
$\binom mq\le\exp(mH(q/m))$, and summing over the $O(\ell)$ permissible
values of $q$, gives, with the higher-omission factor included,
\begin{equation}
 \#\{e_1\}\le
 B^{\frac12 gH(2r_z)+\varepsilon_L(\tau)+o(1)}.
 \label{eq:26}
\end{equation}
The function $\varepsilon_L(\tau)$ can be chosen to tend to zero for
fixed $L$.  To justify this uniformly in the counts, divide $m,q$ by
$\ell$ and use uniform continuity of $H$ on $[0,1]$; the perturbations
are $O(\tau/g)$, and $g\ge g_0>0$ is fixed.  The clipping from $r_z^*$ to
$r_z$ changes these arguments by the same amount.  Also
$\ell/\log B\to1$, and all polynomial factors in $\ell$ contribute
$B^{o(1)}$.  We use $\varepsilon_L(\tau)$ below for a possibly enlarged
function with this same limiting property.

The choices of $d_1$ have the corresponding bound with $r_y$.
Conditional on a regular first remainder $\mathcal T_b$, the number
of its subsets $y$ of the prescribed cardinality is at most
\[
 B^{\frac12 gH(2r_y)+\varepsilon_L(\tau)+o(1)}.
\]
This bound is uniform in that remainder set.  We may condition on it
and sum over these $y$ while the remainder at the other endpoint still
has its independent tilted law.

\paragraph{Availability of a fixed numeric addition.}
The entropy bounds count the possible omitted prime sets. To control
large additions, we also need the probability that each numerically
specified product is present in the other remainder set. The next bound
keeps the first remainder's regularity requirement in this probability.
For a fixed squarefree $z$ in this class, with
$\mathcal P(z)\subset\mathcal Q_g\setminus\mathcal P(a)$, its inclusion
probability in $\mathcal T_a$ is
\[
 \prod_{p\mid z}\frac1{2p-1}
 =\frac{2^{-\omega(z)}}{z}
       \prod_{p\mid z}\left(1-\frac1{2p}\right)^{-1}
 \ll \frac{2^{-\omega(z)}}{z}.
\]
The last bound is uniform because $\log z\le3B$ and all its primes
exceed $P_0$.  Conditional on those inclusions, the remaining available
indicators in $\mathcal Q_g$ are independent with parameter sum
\[
 \Lambda=\sum_{p\in\mathcal Q_g\setminus
                       (\mathcal P(a)\cup\mathcal P(z))}
                       \frac1{2p-1}
         =\tfrac12g\ell+o(\ell).
\]
The excluded primes have reciprocal sum $O(B/P_0)$, so this estimate
is uniform in the fixed coefficient and numeric addition.  First
prefix regularity requires the number of these other inclusions to be
at most $(1/2-r_z^*)g\ell+\tau\ell$.

If $X$ is a sum of independent Bernoulli variables with parameter sum
$\Lambda$, then for $t\ge0$,
\[
 \Prob(X\le k)
 \le\exp\bigl(tk+(e^{-t}-1)\Lambda\bigr).
\]
For $0\le k\le\Lambda$, minimizing gives exponent
$-\Lambda+k-k\log(k/\Lambda)$, with its continuous value at $k=0$.
For $k\ge\Lambda$ use the bound $1$, and for $k<0$ the event is empty.
Uniform continuity of the resulting rate, including the endpoints,
and the fixed-grid relation $g\ge g_0$ therefore imply
\begin{equation}
 \begin{split}
 &\Prob^{\mathrm{tilt}}_{a,b}
   \bigl(\mathcal P(z)\subset\mathcal T_a,
              \ \mathcal T_a\text{ regular at }g\bigr)\\
 &\hspace{12mm}\ll
   \frac{(1/2)^{\omega(z)}}{z}
       B^{-gI_{1/2}(1/2-r_z)+\varepsilon_L(\tau)+o(1)},
 \qquad
 I_{1/2}(u)=\tfrac12-u+u\log\frac{u}{1/2}.
 \end{split}
 \label{eq:27}
\end{equation}
Here $u\log u=0$ at $u=0$.  Retaining only prefix regularity on the
left gives an upper bound for retaining all of first-remainder
regularity, which is what the multiplicity expectation requires.

\paragraph{A harmonic sieve for the numeric addition.}
We now sum the availability bound over numeric additions satisfying
the alternative coefficient relation. The factor $1/z$ in that bound
is the reason for the harmonic normalization below.
Fix $e_1,d_1,y$ from the preceding choices. The alternative size
conditions imply
\[
 z\asymp Z':=d_1y/e_1
\]
with absolute comparison constants; we may enlarge the interval to
$[Z'/2,2Z']$.

For these fixed data and the fixed cardinality class, let
$\mathcal Z$ consist of the positive integers $z$ with the following
properties:
\[
\begin{gathered}
 z\text{ is a squarefree product of primes in }
       \mathcal Q_g\setminus\mathcal P(a),\qquad
 \omega(z)=r_z^*g\ell,\\
 Z'/2\le z\le2Z',\qquad
 \log z\in
 \begin{cases}
 (B^{g-\Delta},B^g],&g<1,\\
 (B^{1-\Delta},3B],&g=1,
 \end{cases}\\
 e_1z\equiv d_1y\pmod{|j|},\qquad
 c'(z):=\frac{e_1z-d_1y}{j}>0,\qquad
 (e_1z,d_1y,c'(z))\in\mathcal C_j .
\end{gathered}
\]
The congruence makes $c'(z)$ an integer. Every valid alternative in
this designated class yields an element of $\mathcal Z$. The inclusion
$\mathcal P(z)\subset\mathcal T_a$ remains the event estimated in
\eqref{eq:27}; the two alternative remainder regularity conditions
are discarded. Equation~\eqref{eq:27} applies to every
$z\in\mathcal Z$, because each is a possible product in the tilted
remainder at the $a$ endpoint. If $\mathcal Z$ is empty there is nothing
to prove; otherwise its interval and log conditions give
$Z'\ge\tfrac12\exp(B^{g-\Delta})$.

All of $e_1,d_1,y$ are units modulo $|j|$, since their prime factors
exceed $P_0>T$. Thus the displayed congruence is one unit residue class.
We claim
\begin{equation}
 \sum_{z\in\mathcal Z}\frac{(1/2)^{\omega(z)}}{z}
 \ll\frac1{|j|}
 B^{g[F_0(r_z)-1-h_0]+3\Delta+r_{\min}+\tau\log2+o(1)},
 \qquad
 F_0(u)=u\left(1+\log\frac{1/2}{u}\right),\quad F_0(0)=0,
 \label{eq:28}
\end{equation}
where $r_{\min}\in(0,1/2)$ is any fixed number. The excess exponent
$3\Delta+r_{\min}+\tau\log2+o(1)$ can be made small by the later choices
of grid, clamp, and tolerance.

To prove the deterministic estimate, put
$\lambda=\max(r_{\min},r_z)$, so $r_{\min}\le\lambda\le1/2$.
On $\mathcal Z$ the fixed cardinality gives
\[
 (1/2)^{\omega(z)}
 =\lambda^{\omega(z)}
       \left(\frac{1/2}{\lambda}\right)^{r_z^*g\ell}.
\]
The membership of $(e_1z,d_1y,c'(z))$ in $\mathcal C_j$ also gives
the upper prefix cutoff for $c'(z)$, and hence
\[
 1\le2^{(g/2+\tau)\ell}\,2^{-\omega_g(c'(z))}.
\]
The prefix factor is introduced to supply a second reducing root in the
sieve below. At the leading $g$-scale its saving $-g/2$ combines with
the cutoff cost $g\log2/2$, leaving $-gh_0$ in the exponent.
The calculation below verifies the root conditions and retains the grid
and tolerance losses.
Insert these two factors before enlarging the summation domain, and
keep only the reducing prime factors up to
\[
 Z_{\mathrm{sieve}}=\exp(B^{g-2\Delta}).
\]
Deleting the other reducing factors increases each summand. Since
$z\ge Z'/2$ on $\mathcal Z$, we obtain the deterministic majorant
\[
\begin{split}
 \sum_{z\in\mathcal Z}\frac{(1/2)^{\omega(z)}}z
 &\le \frac{2}{Z'}
       \left(\frac{1/2}{\lambda}\right)^{r_z^*g\ell}
       2^{(g/2+\tau)\ell}\\
 &\quad\cdot
 \sum_{\substack{z\in\Z\cap[Z'/2,2Z']\\
                 e_1z\equiv d_1y\pmod{|j|}}}
 \prod_{P_0<p\le Z_{\mathrm{sieve}}}
       \lambda^{\ind_{\{p\mid z\}}}
       2^{-\ind_{\{p\mid c'(z)\}}}.
\end{split}
\]
The finite congruence product is defined for every integer value of
$c'(z)$, including zero. The hard prefix cutoff, the fixed-cardinality
condition, and all other candidate restrictions have been removed only
in this deterministic sum. Equation~\eqref{eq:27} is applied only to
$z\in\mathcal Z$; the displayed enlargement bounds the weighted sum
that results.

For explicit congruence accounting put $J_1=|j|$ and
$\sigma=j/J_1$, and choose a unit $z_0$ modulo $J_1$ satisfying the
progression condition.  Write
\[
 z=z_0+J_1t,
 \qquad
 c'=c_0+\sigma e_1t,
 \qquad
 c_0=\frac{e_1z_0-d_1y}{j}\in\Z.
\]
The parameter $t$ runs through an interval of length
$N\asymp Z'/J_1$.  Its logarithm is at least
$B^{g-\Delta}-O(\log B)$.  Consequently
\[
 \frac{\log Z_{\mathrm{sieve}}}{\log N}=O(B^{-\Delta})\longrightarrow0,
\]
so the interval upper sieve of Section~\ref{sec:arithmetic}, in fixed
dimension two, applies uniformly.

The displayed product already omits every prime at most $P_0$, including
every prime of $j$. Among the remaining primes also drop those dividing
$e_1d_1y$.  Their reciprocal sum is $O(B/P_0)$, uniformly, since
$\log(e_1d_1y)=O(B)$.  At each retained prime the two slopes are
invertible, and the two roots are distinct.  Indeed their determinant
is
\[
 J_1c_0-\sigma e_1z_0=-\sigma d_1y,
\]
which is nonzero modulo such a prime.  The local weights of the roots
are $\lambda$ and $1/2$, so the average local weight is exactly
\[
 1-\frac{3/2-\lambda}{p}.
\]
This also explains why no coefficient-height factor enters the sieve:
only the number and distinctness of residue roots are used.  If one
keeps primes dividing $j$ instead, the first form is a fixed unit and
the second has invertible slope; dropping them avoids any need for a
separate local factor.

The random-weight interval sieve bounds the sum of the retained
reducing weights by a constant times
\[
 N\prod_{\substack{P_0<p\le Z_{\mathrm{sieve}}\\p\nmid e_1d_1y}}
       \left(1-\frac{3/2-\lambda}{p}\right)
       +O(N^{.8}).
\]
The first term is
$N B^{(g-2\Delta)(\lambda-3/2)+o(1)}$, by Mertens' estimate.
Indeed the reciprocal sum in that product is
$(g-2\Delta)\log B-\log\log P_0+O(1)
 =(g-2\Delta)\log B+o(\log B)$, and the sum of the squared
reciprocals is negligible.  Discarding the small-prime conditions
costs at most powers of $\log P_0$ relative to their possible sieve
factors, and those powers are $B^{o(1)}$; alternatively, the displayed
large-prime product already proves the upper bound directly.

Using $z\asymp Z'$ and restoring the threshold factors, the harmonic
normalization of the main sieve bound is $N/Z'\asymp1/J_1$.  The
remainder contributes at most
\[
 \frac1{J_1} B^{O(1)}N^{-.2},
\]
which is negligible uniformly in the class, since $g\ge4\Delta$ and
$\log N\ge B^{g-\Delta}-O(\log B)$.  The total power of $B$ in the
main term is
\[
 (g-2\Delta)(\lambda-3/2)
       +g r_z^*\log\frac{1/2}{\lambda}
       +(g/2+\tau)\log2+o(1).
\]
If $r_z^*\le1/2$, then $r_z=r_z^*$, and this power equals
\[
 \begin{split}
 &g[F_0(r_z)-1-h_0]
   +g\left(\lambda-r_z+r_z\log\frac{r_z}{\lambda}\right)\\
 &\hspace{20mm}
   +2\Delta(3/2-\lambda)+\tau\log2+o(1).
 \end{split}
\]
We used $-3/2+(\log2)/2=-1-h_0$.  The clamp contribution in
parentheses vanishes when $r_z\ge r_{\min}$ and lies between $0$ and
$r_{\min}$ when $0\le r_z<r_{\min}$, including the value $r_{\min}$ at
$r_z=0$.  The grid contribution is at most $3\Delta$.
If $r_z^*>1/2$, our choice is $r_z=\lambda=1/2$; the Rankin factor is then
exactly $1$ and the same upper bound follows directly with the clipped
$r_z$.  These observations prove \eqref{eq:28}.  If $y$ is the designated
large numeric variable, the identical proof uses the slope $-d_1$
for $c'$ instead, and its determinant has the same nonvanishing
property outside the already discarded coefficient primes.

\paragraph{Combining the bounds and choosing parameters.}
Condition on the first remainder at the $b$ endpoint.  If it is not
regular its contribution vanishes.  If it is regular, enumerate the
choices of $d_1,y$ by the two entropy bounds with $r_y$, and enumerate
$e_1$ by \eqref{eq:26}.  For each such choice, sum the availability
bound \eqref{eq:27} using \eqref{eq:28}.  Independence of the other
first remainder justifies this conditioning, and every bound is
uniform in the conditioned regular remainder.  Thus we may average
it out without any additional factor.

A direct expansion of $H$ gives the exact identity
\[
 F_0(r_z)-I_{1/2}(1/2-r_z)=\tfrac12H(2r_z)
        \qquad(0\le r_z\le1/2).
\]
The three combinatorial factors have entropy exponents
$gH(2r_z)/2$, $gH(2r_y)/2$, and $gH(2r_y)/2$.  Adding the probability and
sieve exponents therefore bounds the contribution of one class,
one orientation, and one count pair, with both first-remainder
regularity indicators, by
\begin{equation}
 \begin{split}
 &\ll \frac1{|j|}
 B^{g[H(2r_z)+H(2r_y)-1-h_0]
             +3\Delta+r_{\min}+\varepsilon_L(\tau)+o(1)}\\
 &\le \frac1{|j|}
 B^{\kappa_1+3\Delta+r_{\min}
                         +\varepsilon_L(\tau)+o(1)}.
 \end{split}
 \label{eq:29}
\end{equation}
The tolerance error here includes the preceding $\tau\log2$ and the
finitely many entropy and rate perturbations.  The last inequality
uses $H\le\log2$ and $0<g\le1$.  If the bracket in the first exponent
is negative, multiplying it by $g$ still gives a nonpositive number;
otherwise its maximum is at most
$2\log2-1-h_0=\kappa_1$.

Here is an explicit consistent order of choices.  First choose $L$
large enough that $6\Delta\log2<.004$ and $3\Delta<.01$.
Next choose a fixed $r_{\min}<.01$.  Finally make $\tau$ small enough
that the small-addition exponent is strictly below $.008$, that
\[
 3\Delta+r_{\min}+\varepsilon_L(\tau)<.04,
\]
and that all the strict weight inequalities in
\eqref{eq:arith-weight-exponents} hold.
These choices are possible because $\varepsilon_L(\tau)$ and
$\varepsilon_{\mathrm{high}}(\tau)$ tend to zero for the fixed grid.
They impose no condition on $C_*$, since only the prefix and total
regularity conditions were needed for this estimate.

There are finitely many grid classes and orientations and
$O(\ell^2)=B^{o(1)}$ count pairs.  Since
$|j|\ge\eta T\asymp_\eta B^{.32}$ and
\[
 \kappa_1+.04-.32=-.047132048\ldots<0,
\]
the sum of all the large-addition bounds is $o(1)$.
The small-addition bound, with a strict exponent below $.008$, is
$O(B^{.009})$ after the $o(1)$ in its exponent is absorbed.
This proves the lemma, including the first candidate itself when valid.
\end{proof}

\subsection{Squaring the kernel and averaging the rows}

The multiplicity estimate has controlled how many alternative
representations survive after weighting a first one. Combining it with
the small weight of each representation now gives the required
integrated row-square bound.

\begin{proof}[Proof of Proposition~\ref{prop:moments:rows}]
For an allowed edge $i,k=i+j$, apply Lemma~\ref{lem:moments:multiplicity} to
\eqref{eq:moments-first-representation} and use the uniform bound on
$\Xi$. This gives
\[
 \E\mathcal L_{ik}^2
 \ll B^{-\kappa_1+5\tau\log2+.009+o(1)}
 \frac1B\sum_{(a,b,c)\in\mathcal C_j}
       \frac{A(a)A(b)A(c)}{ab}.
\]
The remaining coefficient sum is bounded using
\eqref{eq:7}.  Cover the support of $c$ by $O(B)$ dyadic boxes
$c\asymp X$, so $a,b\asymp TX$ in each box.  On such a box, dropping
coefficient cutoffs for an upper bound gives
\[
 \frac1B\sum_{a-b=jc}\frac{A_0(a)A_0(b)A_0(c)}{ab}
 \ll \frac{1}{B(TX)^2}\,TX^2\Sigma(j)
 =\frac{\Sigma(j)}{BT}.
\]
Summing the boxes yields $O(\Sigma(j)/T)$.  The same proof works for
negative $j$ by symmetry.  Consequently
\[
 \E\mathcal L_{ik}^{2}
 \ll B^{-\kappa_1+5\tau\log2+.009+o(1)}
                  \frac{\Sigma(k-i)}{T}.
\]
The choice in \eqref{eq:arith-weight-exponents} gives
$\kappa_1-5\tau\log2=1-5h_\tau>.229$.
Thus the exponent on the right is strictly below $-.220$ before the
$o(1)$ term.  Finally
\[
 \frac1T\sum_{0<|j|\le T}\Sigma(j)\ll1,
\]
so summing over a row proves \eqref{eq:24}, with room to absorb the
$o(1)$ in the exponent.

For the last assertion set
$D_i=\sum_k\mathcal L_{ik}$.  Given $S_i$, the summands with different
$k\ne i$ are independent, because the other site types are independent
and the roots have already been averaged out.  Therefore
\[
 \E\operatorname{Var}(D_i\mid S_i)
 \le\sum_k\E\mathcal L_{ik}^{2}\ll B^{-.21}.
\]
The conditional mean is bounded for every $S_i$ by \eqref{eq:23}.
The variance identity consequently gives
\[
 \E D_i^2
 =\E\operatorname{Var}(D_i\mid S_i)
       +\E\bigl(\E[D_i\mid S_i]\bigr)^2\ll1.
\]
Notice that only an integrated conditional variance bound has been
used; a pointwise conditional second-moment estimate is unnecessary.
Jensen's inequality across the rows now gives
\[
 \E\left(\frac1M\sum_iD_i\right)^2
 \le\frac1M\sum_i\E D_i^2\ll1.
\]
For the kernel restricted to $I$ as in the statement, nonnegativity can
only reduce the conditional means, row-square sums, and total mass.
Keeping the same normalization $M$ therefore gives the stated bounds.
\end{proof}
